\documentclass[11pt,a4paper]{article}

% ====================== PAKIETY ======================
\usepackage[utf8]{inputenc}
\usepackage[T1]{fontenc}
\usepackage[polish]{babel}
\usepackage[a4paper,margin=2.3cm]{geometry}
\usepackage{lmodern}
\usepackage{graphicx}
\usepackage{array}
\usepackage{longtable}
\usepackage{booktabs}
\usepackage{tabularx}
\usepackage{enumitem}
\usepackage{xcolor}
\usepackage{titlesec}
\usepackage{fancyhdr}
\usepackage{tikz}
\usetikzlibrary{shapes.geometric, shapes.multipart, arrows.meta, positioning, fit, backgrounds, calc}
\usepackage[hidelinks]{hyperref}
\usepackage{caption}
\usepackage{placeins}

% ====================== KOLORY ======================
\definecolor{primary}{RGB}{30,90,150}
\definecolor{accent}{RGB}{0,140,120}
\definecolor{lightgray}{RGB}{240,242,245}
\definecolor{midgray}{RGB}{170,178,186}
\definecolor{riskhigh}{RGB}{200,60,60}
\definecolor{riskmed}{RGB}{220,150,40}
\definecolor{risklow}{RGB}{90,160,90}
\definecolor{passgreen}{RGB}{40,150,70}

% ====================== STYLE NAGŁÓWKÓW ======================
\titleformat{\section}{\normalfont\Large\bfseries\color{primary}}{\thesection.}{0.6em}{}
\titleformat{\subsection}{\normalfont\large\bfseries\color{accent}}{\thesubsection.}{0.6em}{}
\titleformat{\subsubsection}{\normalfont\normalsize\bfseries\color{primary!80!black}}{\thesubsubsection.}{0.5em}{}
\titlespacing*{\section}{0pt}{16pt}{6pt}
\titlespacing*{\subsection}{0pt}{12pt}{4pt}

% ====================== STOPKA / NAGŁÓWEK ======================
\pagestyle{fancy}
\fancyhf{}
\renewcommand{\headrulewidth}{0.4pt}
\renewcommand{\footrulewidth}{0.4pt}
\fancyhead[L]{\small\color{primary}\textbf{RehabSense} -- Dokument Wizji}
\fancyhead[R]{\small Modernizacja Modułu MediaPipe AI}
\fancyfoot[L]{\small Politechnika Łódzka}
\fancyfoot[R]{\small Strona \thepage}

% ====================== POMOCNICZE ======================
\newcolumntype{Y}{>{\raggedright\arraybackslash}X}
\newcommand{\sysname}{\textbf{RehabSense}}
\setlist{itemsep=2pt, topsep=3pt}

% ---- Style diagramów TikZ ----
\tikzset{
  layer/.style={rectangle, rounded corners, draw=primary, thick, fill=lightgray,
                minimum width=3.4cm, minimum height=1cm, align=center, font=\small},
  ext/.style={rectangle, rounded corners, draw=accent, thick, fill=white,
              minimum width=2.8cm, minimum height=0.9cm, align=center, font=\small},
  box/.style={rectangle, rounded corners=3pt, draw=primary, thick, fill=white,
              minimum width=3.0cm, minimum height=1.1cm, align=center, font=\small},
  ai_comp/.style={rectangle, rounded corners=3pt, draw=accent, very thick, fill=accent!10,
                  minimum width=3.2cm, minimum height=1.1cm, align=center, font=\small\bfseries},
  arr/.style={-{Stealth[length=3mm]}, thick, draw=primary!80},
  arr_accent/.style={-{Stealth[length=3mm]}, thick, draw=accent},
}

\begin{document}

\pagenumbering{gobble}

% ====================== STRONA TYTUŁOWA ======================
\begin{titlepage}
\centering
\vspace*{1.5cm}
{\Huge\bfseries\color{primary} DOKUMENT WIZJI\par}
\vspace{0.4cm}
{\LARGE\bfseries System wspomagania rehabilitacji\par}
\vspace{0.3cm}
{\Huge\bfseries\color{accent} RehabSense 2.0\par}
\vspace{0.8cm}
{\large\textsc{Koncepcja Rozwoju i Udoskonalenia Modułu Rehabilitacji AI z Wykorzystaniem Google MediaPipe}\par}
\vspace{1.2cm}
\rule{0.75\textwidth}{1.2pt}\par
\vspace{1.2cm}

{\large\textbf{Instytucja:} Politechnika Łódzka\par}
\vspace{0.3cm}
{\large\textbf{Przedmiot:} Moduł Sumatywny\par}
\vspace{0.3cm}
{\large\textbf{Dokument:} Dokument Wizji Produktu\par}
\vspace{0.3cm}
{\large\textbf{Wersja:} 2.0-draft\par}
\vspace{0.3cm}
{\large\textbf{Data opracowania:} 13 października 2026\par}

\vfill
{\small Dokument definiuje wizję rozwojową systemu RehabSense, koncentrując się na kluczowym zadaniu modernizacji modułu analizy ruchu AI opartego na MediaPipe w celu osiągnięcia najwyższej dokładności, precyzji biometrycznej i skuteczności terapeutycznej.}
\end{titlepage}

\newpage
\pagenumbering{arabic}

% ====================== LISTA CZŁONKÓW ZESPOŁU ======================
\section*{Lista członków zespołu}
\addcontentsline{toc}{section}{Lista członków zespołu}
\begin{table}[ht]
\centering
\begin{tabularx}{0.85\textwidth}{>{\raggedright\arraybackslash\bfseries}p{6.0cm} Y}
\toprule
\textbf{Imię i nazwisko} & \textbf{Rola w projekcie} \\
\midrule
Kacper Błaszczyk (251485) & Project Manager \\
Klim Hudzenko (253833) & Backend Developer \\
Maciej Miazek (251589) & Frontend \& AI Engineer \\
Remigiusz Tomecki (251652) & Frontend Developer \\
Kacper Romek (251619) & Tester \\
\bottomrule
\end{tabularx}
\caption{Skład zespołu projektowego RehabSense.}
\end{table}

\newpage
\tableofcontents
\newpage

% ====================================================
\section{Informacje o dokumencie}
% ====================================================

\subsection{Cel dokumentu}
Celem niniejszego dokumentu jest przedstawienie spójnej i kompleksowej wizji rozwoju produktu \sysname{} w wersji 2.0, jego wartości biznesowej, docelowych grup użytkowników, zakresu funkcjonalnego oraz kluczowych założeń.

Głównym motywem przewodnim nowej iteracji systemu jest \textbf{udoskonalenie modułu rehabilitacji wykorzystującego technologię Google MediaPipe}. Dokument precyzuje, w jaki sposób dotychczasowy moduł wizyjny zostanie zmodernizowany, aby wyeliminować ograniczenia zidentyfikowane w wersji 1.0.

\subsection{Historia zmian}
Poniższa tabela przedstawia historię ewolucji dokumentacji oraz kamieni milowych projektu.

\begin{table}[ht]
\centering
\small
\begin{tabularx}{\textwidth}{p{1.6cm} p{2.2cm} p{3.6cm} Y}
\toprule
\textbf{Wersja} & \textbf{Data} & \textbf{Autorzy} & \textbf{Opis wprowadzonych zmian} \\
\midrule
1.0 & 19.06.2026 & Zespół projektowy & Opracowanie Dokumentu Architektury, Planu Testów, Instrukcji Użytkownika, wdrożenie wersji 1.0 systemu oraz publikacja Raportu z Testów. \\
\textbf{2.0} & \textbf{14.10.2026} & \textbf{Zespół projektowy} & \textbf{Opracowanie kompleksowego Dokumentu Wizji dla RehabSense 2.0; ukierunkowanie rozwoju na udoskonalenie modułu rehabilitacji MediaPipe.} \\
\bottomrule
\end{tabularx}
\caption{Historia zmian dokumentu wizji.}
\end{table}

\subsection{Definicje i skróty}
\begin{itemize}
    \item \textbf{MediaPipe (Google AI Edge):} Zestaw bibliotek uczenia maszynowego pracujących po stronie klienta (przeglądarka), służący do detekcji punktów orientacyjnych ciała w czasie rzeczywistym.
    \item \textbf{DTW (Dynamic Time Warping):} Algorytm porównywania sekwencji czasowych o zróżnicowanej dynamice tempa wykonania.
    \item \textbf{DDTW (Derivative Dynamic Time Warping):} Rozszerzenie algorytmu DTW uwzględniające pierwszą pochodną kątową (prędkość kątową), zapobiegające błędnemu dopasowywaniu pozycji statycznych do faz dynamicznych.
    \item \textbf{JWT (JSON Web Token):} Bezstanowy mechanizm uwierzytelniania użytkowników w komunikacji REST API.
    \item \textbf{NeonDB:} Chmurowa baza danych PostgreSQL wykorzystywana w architekturze backendowej.
\end{itemize}

% ====================================================
\section{Problem biznesowy}
% ====================================================

\subsection{Obecna sytuacja}
Tradycyjny proces rehabilitacji ruchowej opiera się na bezpośrednich wizytach pacjenta w gabinecie fizjoterapeutycznym. Po instruktażu gabinetowym pacjent otrzymuje zalecenie wykonywania serii ćwiczeń w domu. Wdrożona dotychczas wersja 1.0 aplikacji \sysname{} wprowadziła automatyzację tego procesu poprzez webowy interfejs z kamerą internetową, łączący pacjenta ze specjalistą oraz analizę MediaPipe z algorytmem DTW.

\subsubsection{Jak proces działa obecnie?}
\begin{enumerate}
    \item Fizjoterapeuta tworzy plan rehabilitacji i przypisuje ćwiczenia zawierające referencyjne wideo instruktażowe.
    \item Pacjent uruchamia sesję w przeglądarce przed kamerą internetową.
    \item Frontend pobiera wideo referencyjne i odtwarza je obok obrazu z kamery pacjenta.
    \item Moduł analizy uruchamia model detekcji pozy dwukrotnie w każdej klatce: raz dla strumienia wideo z kamery i raz dla klatki wideo referencyjnego.
    \item Z wykrytych punktów 2D obliczane jest 8 płaskich kątów stawowych (łokcie, ramiona, biodra, kolana).
    \item Wyniki poddawane są analizie z wykorzystaniem DTW, wyliczającej procent dopasowania ruchu, a zliczenie powtórzenia następuje w momencie zapętlenia wideo wzorcowego.
\end{enumerate}

\subsubsection{Jakie występują problemy?}
\begin{itemize}
    \item \textbf{Analiza 2D podatna na błędy perspektywy:} Obliczanie kątów w przestrzeni ekranu sprawia, że niewielki obrót pacjenta względem osi kamery zmienia pozorne kąty stawowe, generując fałszywe powiadomienia o błędach.
    \item \textbf{Brak detekcji oszustwa:} Pacjent może oszukiwać test, unosząc całe ciało zamiast zginąć przedramię.
    \item \textbf{Zliczanie powtórzeń zależne od wideo:} Jeśli pacjent ćwiczy wolniej, szybciej lub z przerwami, zliczenia nie są zsynchronizowane.
\end{itemize}

\subsubsection{Jakie są koszty lub ryzyka?}
\begin{itemize}
    \item \textbf{Ryzyko pogłębienia urazu i utrwalenia patologicznych wzorców:} Niepoprawnie wykonywane ćwiczenie bez korekcji kompensacji może prowadzić do przeciążenia sąsiednich stawów.
    \item \textbf{Utrata zaufania fizjoterapeutów:} Terapeuci nie mogą polegać na wskaźnikach dokładności, jeśli są one obarczone artefaktami perspektywy i oświetlenia.
    \item \textbf{Frustracja i spadek motywacji pacjentów:} Niesprawiedliwe obniżanie oceny lub brak zaliczania powtórzeń z powodu czysto technicznych zakłóceń prowadzi do porzucania domowej terapii.
\end{itemize}

\subsection{Problem do rozwiązania}
Zadaniem stojącym przed projektem jest gruntowna modernizacja modułu analizy AI - uodpornienie analizy na zakłócenia i próby oszustwa.

% ====================================================
\section{Wizja produktu}
% ====================================================

Zgodnie z przyjętym standardem inżynierii wymagań, wizja produktu definiuje się następująco:

\vspace{0.3cm}
\noindent\fbox{%
\parbox{0.97\textwidth}{%
\textbf{Dla} pacjentów realizujących domową rehabilitację ruchową,\\
\textbf{którzy potrzebują} specjalistycznego wsparcia ze strony fizjoterapeutów w procesie powrotu do pełnej sprawności,\\
\textbf{produkt} \sysname{} 2.0\\
\textbf{jest} inteligentnym, webowym systemem nadzoru rehabilitacji,\\
\textbf{który umożliwia} adaptacyjną detekcję ruchów, śledzenie postępów wykonywanych ćwiczeń, zapewniając przy tym wysoką dokładność i bezpieczeństwo terapii domowej.
}%
}
\vspace{0.3cm}

% ====================================================
\section{Cele biznesowe}
% ====================================================

\begin{enumerate}
    \item \textbf{Zwiększenie dokładności oceny ruchu:} Podniesienie skuteczności rozpoznawania poprawności ćwiczeń do poziomu co najmniej 85--90\% poprzez lepsze dopasowanie punktów ciała pacjenta.
    \item \textbf{Wykrywanie nieprawidłowości i oszustw:} Wprowadzenie kontroli stabilności tułowia, aby pacjent nie mógł zaliczyć ćwiczenia ruchem całego ciała zamiast ćwiczonego stawu.
    \item \textbf{Poprawne zliczanie powtórzeń:} Zliczanie powtórzeń w oparciu o rzeczywisty ruch pacjenta (faza zgięcia i wyprostu), a nie na podstawie długości filmu wzorcowego.
    \item \textbf{Wiarygodne dane dla fizjoterapeuty:} Zapewnienie, że wyniki i statystyki widoczne w systemie rzetelnie odzwierciedlają jakość wykonanych ćwiczeń.
\end{enumerate}

% ====================================================
\section{Interesariusze}
% ====================================================
Głównymi interesariuszami projektu są:

\begin{itemize}
    \item \textbf{Pacjenci:} Osoby wykonujące ćwiczenia w domu, które oczekują dokładnej oceny swoich ruchów oraz natychmiastowej informacji o błędach.
    \item \textbf{Fizjoterapeuci:} Specjaliści zlecający ćwiczenia, którzy potrzebują pewności, że pacjent wykonuje je poprawnie i bezpiecznie.
    \item \textbf{Zespół projektowy:} Programiści odpowiedzialni za implementację ulepszeń w module analizy MediaPipe i utrzymanie aplikacji.
\end{itemize}

% ====================================================
\section{Grupy użytkowników}
% ====================================================
W systemie występują trzy grupy użytkowników:

\begin{itemize}
    \item \textbf{Pacjent:} Uruchamia sesję ćwiczeń przed kamerą internetową, otrzymuje informację zwrotną o poprawności ruchu i widzi liczbę poprawnie wykonanych powtórzeń.
    \item \textbf{Fizjoterapeuta:} Tworzy plany rehabilitacji, przegląda historię pacjenta i sprawdza dokładność wykonania zadanych ćwiczeń.
    \item \textbf{Administrator:} Zarządza kontami w systemie i weryfikuje profile fizjoterapeutów.
\end{itemize}

% ====================================================
\section{Zakres produktu}
% ====================================================

\subsection{Funkcjonalności w zakresie}
Prace rozwojowe skupiają się wyłącznie na module analizy wykonywanych ćwiczeń:

\begin{enumerate}
    \item \textbf{Dopasowanie punktów do sylwetki:}
    Przeliczanie współrzędnych punktów z MediaPipe względem wielkości tułowia pacjenta, dzięki czemu drobne przesunięcia lub lekki obrót względem kamery nie zafałszowują wyników.
    
    \item \textbf{Wykrywanie potencjalnych oszustw:}
    Sprawdzanie, czy tułów pacjenta pozostaje stabilny podczas ruchu kończyny. Jeśli pacjent zamiast zgiąć rękę lub nogę przechyla całe ciało, system nie zalicza powtórzenia i informuje o błędzie.
    
    \item \textbf{Licznik powtórzeń oparty na fazach ruchu:}
    Automatyczne rozpoznawanie momentu rozpoczęcia ruchu, osiągnięcia pełnego zgięcia oraz powrotu do pozycji wyjściowej. Zliczanie powtórzeń staje się niezależne od tempa odtwarzania wideo wzorcowego.
    
    \item \textbf{Ulepszone porównywanie ruchu:}
    Większy wpływ na ocenę mają stawy kluczowe dla danego ćwiczenia (np. łokieć przy zginaniu ramienia), a mniejszy stawy, które powinny pozostać nieruchome.
    
    \item \textbf{Prosty komunikat o błędzie na ekranie:}
    Wyświetlanie czytelnego ostrzeżenia wizualnego w momencie, gdy pacjent wykonuje ruch niepoprawnie (np. „Ustabilizuj tułów”).
\end{enumerate}

\subsection{Poza zakresem}
Z zakresu prac wyłączono:
\begin{itemize}
    \item Przejście na modele trójwymiarowe (3D) oraz integrację z dodatkowymi czujnikami zewnętrznymi.
    \item Modyfikację innych modułów aplikacji (np. czat, rejestracja, system logowania pozostają bez zmian).
    \item Automatyczne diagnozowanie schorzeń przez sztuczną inteligencję (system ocenia jedynie technikę ćwiczeń zleconych przez fizjoterapeutę).
\end{itemize}

% ====================================================
\section{Kluczowe wymagania biznesowe}
% ====================================================
\begin{table}[ht]
\centering
\small
\begin{tabularx}{\textwidth}{l p{4.5cm} Y c}
\toprule
\textbf{ID} & \textbf{Wymaganie} & \textbf{Opis} & \textbf{Priorytet} \\
\midrule
\textbf{1} & Normalizacja pozycji ciała & System przelicza punkty względem tułowia, aby odległość od kamery nie psuła wyniku pomiaru. & Wysoki \\
\textbf{2} & Blokada potencjalnych oszustw & System wykrywa ruch tułowia przy ćwiczeniach kończyn i nie zalicza takiego powtórzenia. & Wysoki \\
\textbf{3} & Niezależny licznik powtórzeń & System zlicza powtórzenia po wykonaniu wszystkich faz ćwiczenia & Wysoki \\
\textbf{4} & Ocena DTW & Większą wagę przy ocenie mają stawy wykonujące ruch w danym ćwiczeniu. & Wysoki \\
\textbf{5} & Komunikat o błędzie techniki & Wyświetlanie ostrzeżenia na ekranie, gdy ćwiczenie wykonywane jest niepoprawnie. & Średni \\
\bottomrule
\end{tabularx}
\caption{Kluczowe wymagania biznesowe.}
\end{table}

% ====================================================
\section{Główne procesy biznesowe}
% ====================================================

\subsection{Proces 1: Rozpoczęcie sesji ćwiczeń}
\begin{enumerate}
    \item Pacjent wybiera ćwiczenie z planu rehabilitacji i uruchamia kamerę.
    \item System wykrywa sylwetkę pacjenta przy użyciu MediaPipe i ustala punkty odniesienia.
    \item Pacjent przyjmuje pozycję wyjściową i rozpoczyna ćwiczenie.
\end{enumerate}

\subsection{Proces 2: Analiza ruchu i zliczanie powtórzeń}
\begin{enumerate}
    \item Pacjent wykonuje ruch; system na bieżąco oblicza kąty stawów.
    \item System sprawdza, czy tułów pacjenta jest nieruchomy.
    \item Po osiągnięciu odpowiedniego kąta i powrocie do pozycji początkowej licznik powtórzeń jest inkrementowany.
    \item Jeśli pacjent porusza tułowiem, system wyświetla ostrzeżenie i nie zalicza powtórzenia.
    \item Po wykonaniu zadanej liczby powtórzeń seria zostaje zakończona.
\end{enumerate}

\subsection{Proces 3: Zapis wyników i weryfikacja przez fizjoterapeutę}
\begin{enumerate}
    \item Wyniki sesji są przesyłane do bazy danych.
    \item Fizjoterapeuta w swoim panelu widzi podsumowanie i może ocenić, czy pacjent ćwiczył poprawnie.
\end{enumerate}

% ====================================================
\section{Wymagania niefunkcjonalne}
% ====================================================
\begin{enumerate}
    \item \textbf{Dokładność:} Dokładność oceny poprawności ćwiczenia na poziomie minimum 85\%.
    \item \textbf{Wydajność:} Płynna analiza w przeglądarce internetowej (przynajmniej 25--30 klatek na sekundę).
    \item \textbf{Prywatność danych:} Obraz z kamery jest analizowany wyłącznie w przeglądarce pacjenta i nie jest przesyłany na serwer.
    \item \textbf{Prostota obsługi:} Brak konieczności instalowania dodatkowego oprogramowania ani kupowania drogiego sprzętu (wystarczy zwykła kamera internetowa).
\end{enumerate}

% ====================================================
\section{Integracje}
% ====================================================
Moduł analizy korzysta z:
\begin{itemize}
    \item \textbf{Google MediaPipe:} Biblioteka działająca w przeglądarce pacjenta, wykrywająca punkty ciała na obrazie z kamery.
    \item \textbf{Backend (FastAPI) i Baza Danych (NeonDB):} Zapisywanie informacji o wykonanych seriach, liczbie powtórzeń i średniej dokładności.
\end{itemize}

% ====================================================
\section{Założenia i ograniczenia}
% ====================================================
\begin{itemize}
    \item \textbf{Założenia:} Pacjent posiada sprawną kamerkę internetową i wykonuje ćwiczenia w oświetlonym pomieszczeniu w odległości około 1.5--2 metrów od obiektywu.
    \item \textbf{Ograniczenia:} Zwykła kamera 2D widzi ruch tylko w jednej płaszczyźnie, dlatego pacjent powinien ustawiać się przodem lub pod lekkim kątem do kamery zgodnie z instrukcją ćwiczenia.
\end{itemize}

% ====================================================
\section{Ryzyka}
% ====================================================
\begin{table}[ht]
\centering
\small
\begin{tabularx}{\textwidth}{p{4cm} p{4.5cm} Y}
\toprule
\textbf{Ryzyko} & \textbf{Skutek} & \textbf{Działanie zapobiegawcze} \\
\midrule
Słabe oświetlenie w pomieszczeniu & Gorsze wykrywanie punktów ciała przez MediaPipe & Wyświetlenie prośby o poprawę oświetlenia w przypadku problemów z detekcją. \\
\addlinespace
Oszukiwanie ćwiczenia ruchem całego ciała & Niesłuszne zaliczenie powtórzenia & Algorytm sprawdzający bezruch tułowia podczas ćwiczenia rąk lub nóg. \\
\addlinespace
Nietypowe tempo ćwiczenia pacjenta & Błędne zliczanie powtórzeń & Licznik oparty na osiągnięciu kąta skrajnego, a nie na czasie trwania filmu. \\
\bottomrule
\end{tabularx}
\caption{Główne ryzyka i działania zapobiegawcze.}
\end{table}

% ====================================================
\section{Strategia sukcesu}
% ====================================================
Sukces wdrożenia zmian zostanie potwierdzony, gdy:
\begin{itemize}
    \item Błędy zliczania powtórzeń spadną poniżej 5\% przy poprawnym wykonywaniu ćwiczenia.
    \item Próby oszukiwania ruchem tułowia będą wykrywane w co najmniej 80\% przypadków.
    \item Aplikacja będzie działać płynnie bez zacinania obrazu na standardowym laptopie.
\end{itemize}

% ====================================================
\section{Roadmapa produktu}
% ====================================================
Plan realizacji prac nad poprawą modułu analizy:
\begin{enumerate}
    \item \textbf{Krok 1:} Wdrożenie przeliczania punktów względem tułowia oraz wag w algorytmie DTW.
    \item \textbf{Krok 2:} Implementacja sprawdzania stabilności tułowia w celu wykrywania oszustw.
    \item \textbf{Krok 3:} Wdrożenie niezależnego licznika powtórzeń opartego na zgięciu i wyproście.
    \item \textbf{Krok 4:} Testy dokładności z udziałem użytkowników i dopracowanie progów czułości.
\end{enumerate}

% ====================================================
\section{Korzyści biznesowe}
% ====================================================
\begin{itemize}
    \item \textbf{Dla pacjenta:} Pewność, że ćwiczy poprawnie i bezpiecznie, bez niesłusznego odrzucania poprawnych powtórzeń.
    \item \textbf{Dla fizjoterapeuty:} Rzetelne statystyki pokazujące rzeczywisty wysiłek pacjenta, a nie ruchy zastępcze czy oszukane powtórzenia.
    \item \textbf{Dla projektu:} Rozwiązanie głównego problemu z jakością analizy przy minimalnej złożoności kodu.
\end{itemize}

% ====================================================
\section{Załączniki}
% ====================================================

\subsection{Załącznik 1: Schemat działania ulepszonego modułu analizy}
\begin{figure}[ht]
\centering
\resizebox{0.9\textwidth}{!}{%
\begin{tikzpicture}[node distance=1.2cm and 1.5cm]
  \node[box] (cam) {\textbf{Obraz z kamery}\\ \footnotesize 30 FPS};
  \node[ai_comp, right=1.2cm of cam] (mp) {\textbf{Google MediaPipe}\\ \footnotesize Wykrycie punktów ciała};
  \node[box, right=1.2cm of mp] (norm) {\textbf{Normalizacja punktów}\\ \footnotesize Dopasowanie do tułowia};
  
  \node[box, below=1.2cm of norm] (guard) {\textbf{Sprawdzenie tułowia}\\ \footnotesize Czy pacjent nie oszukuje?};
  \node[box, left=1.2cm of guard] (calc) {\textbf{Pomiar kątów stawów}\\ \footnotesize np. zgięcie łokcia};
  \node[ai_comp, left=1.2cm of calc] (rep) {\textbf{Licznik powtórzeń}\\ \footnotesize Zgięcie i wyprost};

  \node[layer, below=1.2cm of calc, minimum width=8cm] (out) {\textbf{Informacja na ekranie}\\ \footnotesize Aktualny wynik poprawności, liczba powtórzeń i ostrzeżenia};

  \draw[arr] (cam) -- (mp);
  \draw[arr] (mp) -- (norm);
  \draw[arr] (norm) -- (guard);
  \draw[arr] (guard) -- (calc);
  \draw[arr] (calc) -- (rep);
  \draw[arr] (calc) -- (out);
  \draw[arr] (rep) -- (out);
  \draw[arr] (guard) -- (out);
\end{tikzpicture}}
\caption{Przepływ danych w ulepszonym module analizy ćwiczeń.}
\end{figure}

\subsection{Załącznik 2: Kluczowe punkty wykorzystywane do oceny}
\begin{itemize}
    \item \textbf{Tułów (barki i biodra):} Służą jako baza odniesienia do sprawdzania, czy pacjent trzyma prostą postawę i nie buja ciałem.
    \item \textbf{Łokcie i ramiona:} Mierzenie kątów przy ćwiczeniach rąk (zgięcie i wyprost).
    \item \textbf{Kolana i biodra:} Mierzenie kątów przy ćwiczeniach nóg (np. przysiady, unoszenie nóg).
\end{itemize}

\end{document}
